% =====================================================================
%  Glance ML Internship Assignment - Multimodal Fashion & Context Retrieval
%  Compile with:  pdflatex writeup.tex   (run twice for the TOC-free refs)
%  Only standard packages - drops straight into Overleaf.
% =====================================================================
\documentclass[10pt,a4paper]{article}

\usepackage[utf8]{inputenc}
\usepackage[T1]{fontenc}
\usepackage[margin=14mm]{geometry}
\usepackage{booktabs}
\usepackage{tabularx}
\usepackage{array}
\usepackage{xcolor}
\usepackage{enumitem}
\usepackage{titlesec}
\usepackage{parskip}

% Clickable links (GitHub URL, section refs). Loaded last, as hyperref prefers.
\usepackage[colorlinks=true,
            linkcolor=NavyLink,
            urlcolor=NavyLink,
            citecolor=NavyLink,
            pdftitle={Multimodal Fashion and Context Retrieval},
            pdfauthor={Glance ML Internship Assignment}]{hyperref}

\definecolor{NavyLink}{HTML}{1A63C4}
\definecolor{HeadNavy}{HTML}{0B1F3A}
\definecolor{RuleGrey}{HTML}{D7DDE5}

% Compact, tidy headings
\titleformat{\section}{\normalfont\large\bfseries\color{HeadNavy}}{\thesection}{0.6em}{}
\titlespacing*{\section}{0pt}{8pt}{3pt}
\titleformat{\subsection}{\normalfont\normalsize\bfseries\color{HeadNavy}}{\thesubsection}{0.6em}{}
\titlespacing*{\subsection}{0pt}{6pt}{2pt}

\setlist[itemize]{leftmargin=1.2em, topsep=2pt, itemsep=1pt, parsep=0pt}
\setlist[enumerate]{leftmargin=1.4em, topsep=2pt, itemsep=1pt, parsep=0pt}
\renewcommand{\arraystretch}{1.12}
\setlength{\tabcolsep}{4pt}

% Shorthands
\newcommand{\tcode}[1]{\texttt{\small #1}}
\newcommand{\swWin}{\textbf{Y}}
\newcommand{\swTie}{=}
\newcommand{\swLoss}{\textbf{N}}
\newcommand{\repourl}{https://github.com/prakhar-raj-031104/FASHIONATTIRE}

\begin{document}
\pagestyle{plain}

% ---------------------------------------------------------------- title
\begin{center}
  {\LARGE\bfseries\color{HeadNavy} Multimodal Fashion \& Context Retrieval}\\[3pt]
  {\normalsize Glance ML Internship Assignment}\\[4pt]
  {\normalsize \textbf{Codebase:} \href{\repourl}{\repourl}}
\end{center}
\vspace{2pt}
\hrule height 0.8pt
\vspace{6pt}

A natural-language image search engine over \textbf{3{,}200 fashion images} that reasons jointly
about \textbf{what} a person wears, \textbf{where} they are, and the outfit's \textbf{vibe}.

\textbf{Core thesis:} a single global CLIP embedding is \emph{structurally} incapable of
compositional retrieval, so the fix must add \textbf{structure}, not a bigger encoder.

\begin{center}\small
\begin{tabularx}{\textwidth}{@{}>{\raggedright\arraybackslash}p{0.20\textwidth}
                                >{\raggedright\arraybackslash}p{0.14\textwidth}
                                >{\raggedright\arraybackslash}X
                                >{\raggedright\arraybackslash}p{0.12\textwidth}@{}}
\toprule
\textbf{Corpus} & \textbf{Latency} & \textbf{Key measured result} & \textbf{Training} \\
\midrule
3{,}200 images, 9{,}851 garment regions, FAISS $\times$3 + SQLite
& $\sim$150\,ms warm
& Colour-swap overlap@10: vanilla CLIP \textbf{0.750} $\rightarrow$ this system \textbf{0.525}
  (\textbf{$-$30\%}, wins 8/12 pairs)
& \textbf{None} --- fully zero-shot \\
\bottomrule
\end{tabularx}
\end{center}

% ---------------------------------------------------------------- 1
\section{Why this is hard}

\begin{center}\small
\begin{tabularx}{\textwidth}{@{}c>{\raggedright\arraybackslash}p{0.34\textwidth}
                                >{\raggedright\arraybackslash}p{0.20\textwidth}
                                >{\raggedright\arraybackslash}X@{}}
\toprule
\# & \textbf{Query} & \textbf{Capability} & \textbf{CLIP's problem} \\
\midrule
1 & bright yellow raincoat & colour + garment & sparse fine-grained fashion vocabulary \\
2 & business attire in a modern office & style + scene & needs scene grounding \\
3 & blue shirt on a park bench & colour + garment + scene & multi-constraint conjunction \\
4 & casual weekend outfit for a city walk & \textbf{style inference} & \textbf{no garment words at all} \\
5 & a red tie \textbf{and} a white shirt & \textbf{compositional binding} & \textbf{canonical CLIP failure} \\
\bottomrule
\end{tabularx}
\end{center}

\textbf{The actual failure mechanism.} CLIP is trained contrastively over a \textbf{single pooled
embedding per image}. Pooling is order- and position-agnostic: the vector encodes \emph{``there is
red, white, a tie, a shirt''} but \textbf{not which attribute binds to which object}. So
\emph{``red tie + white shirt''} and \emph{``white tie + red shirt''} map to nearly the same point
--- literally the same bag of concepts. This is \textbf{architectural, not capacity}: scaling or
fine-tuning the encoder cannot recover information destroyed at pooling. \textbf{Therefore the fix
must reintroduce structure below the image level.}

% ---------------------------------------------------------------- 2
\section{Approaches and trade-offs \textnormal{\small(Deliverable 1)}}

\begin{center}\small
\begin{tabularx}{\textwidth}{@{}>{\raggedright\arraybackslash}p{0.17\textwidth}
                                >{\raggedright\arraybackslash}X
                                >{\raggedright\arraybackslash}X
                                >{\raggedright\arraybackslash}p{0.19\textwidth}@{}}
\toprule
\textbf{Approach} & \textbf{Strengths} & \textbf{Weaknesses} & \textbf{Best when} \\
\midrule
\textbf{A.} Vanilla CLIP + FAISS & trivial, strong zero-shot & \textbf{bag-of-words}: no binding; weak fashion terms & quick baseline, single-concept queries \\
\textbf{B.} FashionCLIP + FAISS & large domain gain on garment/colour/fabric for $\sim$zero cost & same architectural binding limit & any fashion task --- the \emph{floor}, not the solution \\
\textbf{C.} Caption $\rightarrow$ LLM attributes & interpretable; good scene/style & errors compound (hallucinated caption $\rightarrow$ wrong attribute); non-deterministic; extra GPU & open-vocabulary attributes with reliable captions \\
\textbf{D.} CLIP zero-shot tagging & grounded in \emph{pixels}; deterministic; cheap; editable vocabulary & image-level only $\rightarrow$ still no binding & structured, filterable attributes without training \\
\textbf{E.} Region decomposition & genuinely \textbf{solves binding}; dataset-agnostic & segmentation cost; fails on tiny/occluded items & whenever compositionality is graded \\
\textbf{F.} Detector + scene graph + VQA, or Pinecone/microservices & maximal expressiveness & large complexity, more failure surface, \textbf{no measured gain at this scale} & large production systems \\
\bottomrule
\end{tabularx}
\end{center}

\textbf{Decision: B + D + E}, plus a caption signal from C (BLIP only, \textbf{no LLM extraction}
--- grounding attributes in pixels beats grounding them in possibly-hallucinated text).
\textbf{F rejected} as over-engineering: the brief says favour ML logic over indexing engineering,
so I used FAISS + SQLite and spent the effort on ranking.

% ---------------------------------------------------------------- 3
\section{Chosen architecture \textnormal{\small(Deliverable 2)}}

\textbf{Part A --- Indexer} (\tcode{src/indexing/}). Per image: (1) \textbf{FashionCLIP} global
embedding $\rightarrow$ \tcode{global.faiss}; (2) \textbf{BLIP caption $\rightarrow$ MPNet} sentence
embedding $\rightarrow$ \tcode{caption.faiss}; (3) \textbf{CLIP zero-shot attributes}
(environment/style/garment/colour, each value embedded through several \textbf{prompt templates and
averaged}) $\rightarrow$ SQLite; (4) \textbf{SegFormer garment regions}, each cropped, embedded, and
tagged with \textbf{its own colour and type} $\rightarrow$ \tcode{region.faiss} + SQLite. Artefact (4)
is the crux: it preserves \emph{which colour sits on which garment}.

\textbf{Part B --- Retriever} (\tcode{src/retrieval/}). \emph{Stage 1 (recall):} query embedded by
FashionCLIP + MPNet; union of top-N from the global and caption indexes forms a 400-candidate pool.
\emph{Stage 2 (rerank):} only that pool is scored with expensive signals. This ANN$\rightarrow$rerank
split is why the logic is unchanged at 1M images (\S6).

\begin{center}
\tcode{score = 0.40\,$\cdot$\,global\_clip + 0.20\,$\cdot$\,caption\_sim + 0.20\,$\cdot$\,attribute\_match + 0.20\,$\cdot$\,region\_binding}
\end{center}

\begin{center}\small
\begin{tabularx}{\textwidth}{@{}>{\raggedright\arraybackslash}p{0.19\textwidth}Xc@{}}
\toprule
\textbf{Signal} & \textbf{Fixes} & \textbf{Rescues} \\
\midrule
\tcode{global\_clip} & overall semantics + scene & all \\
\tcode{caption\_sim} & \textbf{style inference} --- caption \emph{``a man in a hoodie walking down a street''} matches \emph{``casual weekend outfit for a city walk''} with \textbf{zero shared garment words} & \textbf{4} \\
\tcode{attribute\_match} & grounded multi-attribute filtering & \textbf{1, 2} \\
\tcode{region\_binding} & \textbf{compositional colour$\leftrightarrow$garment binding} & \textbf{5} \\
\bottomrule
\end{tabularx}
\end{center}

\textbf{Query parsing.} A deterministic rule + synonym + adjacency parser yields a
\tcode{QuerySpec}: structured attributes plus \tcode{(colour, garment)} \textbf{bindings}
(\emph{``a red tie and a white shirt''} $\rightarrow$ \tcode{[(red,tie), (white,shirt)]}). Rules over
an LLM because the schema is small and fixed, it is instant, reproducible and unit-testable ---
while remaining a drop-in point for an LLM parser (\S7.2).

\textbf{Binding score.} A binding scores high only if \textbf{one single region} matches
\emph{both} garment type \emph{and} colour. Multiple bindings combine as
\tcode{0.5$\cdot$mean + 0.5$\cdot$min}: the query is a \textbf{conjunction}, so half-satisfaction
must not score like full (plain \tcode{mean} is too lenient), but plain \tcode{min} is brittle ---
one missed segmentation would zero a perfect image.

\textbf{Two subtleties, both found by measurement:}
\begin{enumerate}
  \item \textbf{Per-signal normalisation is mandatory.} CLIP cosines occupy $\sim$0.15--0.35 while
  attribute/region scores span $[0,1]$; a raw weighted sum would drown the CLIP signal. Each signal
  is min-max normalised \textbf{across the candidate pool} before weighting; signals that do not
  apply are dropped and their weight \textbf{redistributed}.
  \item \textbf{Bound attributes must be excluded from} \tcode{attribute\_match}. That signal is
  \emph{order-agnostic}: both swap variants reduce to $\{$red, white$\}$, so counting them again
  pulls swapped queries back \textbf{together}, diluting the order-sensitive binding. Excluding them
  improved measured separation \textbf{15.8\% $\rightarrow$ 23.7\%}.
\end{enumerate}

% ---------------------------------------------------------------- 4
\section{Measured results}

\subsection{Compositionality --- objective and label-free}

For a query \textbf{Q} and its \textbf{colour-swapped twin Q$'$}, a bag-of-words retriever cannot
tell them apart $\rightarrow$ top-$k$ sets nearly \textbf{identical}; a compositional one returns
\textbf{different} images. So \textbf{overlap@$k$ is an inverse proxy for compositional sensitivity,
needing no human labels.} Both systems run on the \textbf{same index}; only ranking differs.
\textbf{12 pairs $\times$ $k$=10 = 120 slots}, sized so one image (0.008) cannot dominate the mean.

\begin{center}\small
\begin{tabularx}{\textwidth}{@{}Xccp{0.5cm}@{\hspace{4pt}}Xcc@{}}
\toprule
\textbf{Pair} & \textbf{CLIP} & \textbf{Ours} & & \textbf{Pair} & \textbf{CLIP} & \textbf{Ours} \\
\midrule
green top / yellow skirt      & 1.00 & \textbf{0.10}\,\swWin & & blue shirt / red trousers     & 0.90 & \textbf{0.70}\,\swWin \\
purple top / black trousers   & 0.70 & \textbf{0.20}\,\swWin & & black dress / white jacket    & 0.90 & \textbf{0.80}\,\swWin \\
yellow top / blue skirt       & 0.80 & \textbf{0.30}\,\swWin & & black jacket / white trousers & 0.70 & 0.70\,\swTie \\
green jacket / white shirt    & 0.80 & \textbf{0.40}\,\swWin & & pink blouse / black skirt     & 0.70 & 0.70\,\swTie \\
brown coat / black trousers   & 0.90 & \textbf{0.50}\,\swWin & & white blouse / blue jeans     & 0.40 & 0.50\,\swLoss \\
red skirt / white blouse      & 0.70 & \textbf{0.60}\,\swWin & & red tie / white shirt         & 0.50 & 0.80\,\swLoss \\
\midrule
\textbf{MEAN} & \textbf{0.750} & \textbf{0.525} & & \textbf{relative reduction} & \multicolumn{2}{c}{\textbf{$-$30.0\%}} \\
\bottomrule
\end{tabularx}
\end{center}

\textbf{8/12 wins, 2 ties, 2 losses.} Clearest case: \emph{green top / yellow skirt}, where CLIP
returns an \textbf{identical} top-10 for both phrasings (1.00) --- the textbook failure --- while
this system separates them almost completely (0.10).

\textbf{Where gains concentrate.} All 8 wins are pairs where CLIP is most blind (0.70--1.00); both
losses are pairs CLIP \emph{already} separated (0.40, 0.50). \textbf{The mechanism contributes most
precisely where the baseline fails hardest} --- the desired behaviour for a corrective signal, and
it suggests making the binding weight \emph{adaptive} to how confidently the dense signals already
discriminate.

\subsection{Honest negative result, diagnosed}

\emph{Red tie / white shirt} --- query 5 --- \textbf{did not improve}. Traced in data: the index
held only \textbf{18 tie regions / 3{,}200 images}, because two filters removed small garments
(0.5\% area threshold; a 6-region cap applied to an \emph{area-sorted} list, so a tie --- the
smallest item --- was truncated). I lowered the threshold to 0.15\% and raised the cap: regions
\textbf{7{,}321 $\rightarrow$ 9{,}851}, ties \textbf{18 $\rightarrow$ 44} --- and query 5
\emph{still} did not improve.

\textbf{The remaining ceiling is the corpus, not the algorithm.} Fashionpedia is
womenswear/runway-heavy: 44 ties ($\sim$1.4\%), few red. The mechanism demonstrably works where
garments exist (1.00 $\rightarrow$ 0.10 above). The fix is \textbf{dataset composition} (blend in
menswear/formal imagery), not a ranker change. \emph{Two plausible assumptions --- ``regions exist''
and ``more regions will fix it'' --- were both false; only measurement revealed it.}

\subsection{Official queries and corpus}

\begin{center}\small
\begin{tabularx}{\textwidth}{@{}cc>{\raggedright\arraybackslash}p{0.36\textwidth}X@{}}
\toprule
\# & \textbf{Score} & \textbf{Top result} & \textbf{Assessment} \\
\midrule
1 & 0.942 & \emph{``a little girl in a \textbf{yellow raincoat} and red tights''} & exact match \\
2 & 0.916 & \emph{``a woman in a black suit and \textbf{white shirt}''} & attire right; office \emph{scene} is the weak axis \\
3 & 0.906 & \emph{``a young girl sitting on a \textbf{bench in the park}''} & scene + pose correct \\
4 & 0.896 & \emph{``on the \textbf{sidewalk} in a white shirt and \textbf{jeans}''} & inferred from \textbf{no garment words} \\
5 & 0.702 & \emph{``a man in a vest and red pants''} & weakest --- corpus-limited (\S4.2) \\
\bottomrule
\end{tabularx}
\end{center}

Latency \textbf{318\,ms} mean ($\sim$150\,ms warm). Zero-shot environment tagging confirms the
required axes: urban street \textbf{765}, home interior \textbf{391}, park \textbf{292}, office
\textbf{198}; remainder catalogue/runway ``studio'' (1{,}177) --- the Fashionpedia bias, and why
query 2's scene component is weakest.

% ---------------------------------------------------------------- 5
\section{Shortcomings and fixes}

\begin{center}\small
\begin{tabularx}{\textwidth}{@{}>{\raggedright\arraybackslash}p{0.24\textwidth}XX@{}}
\toprule
\textbf{Shortcoming} & \textbf{Cause} & \textbf{Fix} \\
\midrule
Small garments under-detected & area threshold + region cap; only 44 ties/3{,}200 & lower threshold (done); accessory-specialised detector; weight by \emph{saliency} not area \\
Multi-person images bind wrongly & SegFormer parses clothing, not \emph{people} --- a red tie on A and white shirt on B satisfies both & person detector; require bindings to resolve \textbf{within one individual} \\
Uncalibrated attribute confidences & softmax gives \emph{relative}, not probabilistic, scores; scales differ per axis & per-axis Platt/temperature calibration \\
Scene axis corpus-limited & Fashionpedia is catalogue-heavy ($\sim$6\% office) & blend a scene-rich people dataset --- loader is dataset-agnostic by design \\
Parser misses negation/comparatives & ``no tie'', ``more formal than'' out of grammar & small instruction-tuned LLM emitting the \textbf{same} \tcode{QuerySpec} JSON --- interface exists \\
Hand-set weights & no labelled set in the time budget & learn the fusion (\S7.2) \\
\bottomrule
\end{tabularx}
\end{center}

% ---------------------------------------------------------------- 6
\section{Modular code, scalability, zero-shot \textnormal{\small(Criteria 2--4)}}

\textbf{Modular --- logic separated from data.} \tcode{configs/*.yaml} holds every knob (no constant
in code); \tcode{src/models/} wrappers are the \textbf{only} code importing HuggingFace, so
FashionCLIP $\rightarrow$ SigLIP is a one-line YAML change; \textbf{vectors in FAISS, metadata in
SQLite --- never in filenames}; typed dataclasses (\tcode{ImageRecord}/\tcode{RegionRecord}/%
\tcode{QuerySpec}) pass between independently testable stages; the loader assumes only ``a folder of
images'', so Fashionpedia/DeepFashion/COCO/phone photos all work unchanged. \emph{Filename-%
independence is structural: the dataset's filenames are MD5 hashes with zero semantic content, and
\tcode{image\_path} never enters any scoring computation.}

\textbf{Scalability to 1M+.} The retrieval \textbf{logic is size-invariant}; only the index type
changes. \emph{Recall:} \tcode{IndexFlatIP} $\rightarrow$ \tcode{IndexIVFFlat}/HNSW is a
\textbf{one-line} change in \tcode{VectorStore.build(index\_type="ivf")} --- ids and the
\tcode{search()} contract are identical, so nothing else moves; $O(N) \rightarrow \sim O(\sqrt{N})$
or $O(\log N)$. \emph{Rerank:} cost is \textbf{independent of corpus size} --- only the fixed
400-candidate pool is scored --- the key scalability property. \emph{Memory:} 1M $\times$ 512-d fp16
$\approx$ \textbf{1\,GB} global (regions $\sim$2.3$\times$); IVF-PQ compresses 8--32$\times$.
\emph{Metadata:} SQLite is fine into low millions, and structured attributes enable \textbf{SQL
pre-filtering} (\tcode{environment=office AND EXISTS a tie region}) to shrink the ANN space
\emph{before} search. \emph{Indexing:} embarrassingly parallel --- shard and merge FAISS shards.
Measured \textbf{3{,}200 images in $\sim$11 min} on one 4\,GB laptop GPU ($\sim$0.2\,s/image).

\textbf{Zero-shot.} \textbf{Nothing is trained on this corpus or these queries} --- FashionCLIP,
BLIP, MPNet, SegFormer are all off-the-shelf, and attributes come from \textbf{zero-shot CLIP
classification} against an editable vocabulary. A new colour, garment, city or weather condition is
added by editing \tcode{vocab.py} with \textbf{no retraining}. Query 4 is the demonstration:
\emph{``casual weekend outfit for a city walk''} has \textbf{no garment label}, yet returns
sidewalk/jeans imagery because caption and attribute signals infer the vibe rather than matching a
keyword.

% ---------------------------------------------------------------- 7
\section{Future work \textnormal{\small(Deliverable 4)}}

\subsection{Adding locations (cities, places) and weather}

Environment is already a \textbf{first-class structured axis}, so this is mostly \emph{data} plus
new scoring terms --- no architectural rewrite.

\begin{enumerate}
  \item \textbf{Vocabulary extension (no code-path change):} add
  \tcode{weather = [sunny, rainy, snowy, overcast\ldots]} and richer \tcode{place} values to
  \tcode{vocab.py} with prompt templates. FashionCLIP tags them \textbf{zero-shot} immediately; the
  parser gains synonyms. This alone makes \emph{``a raincoat on a rainy street''} prefer genuinely
  wet scenes.
  \item \textbf{Specialised classifiers for precision:} \textbf{Places365} for scene/venue and a
  sky/weather classifier, stored as two more attribute axes and two more weighted terms.
  \item \textbf{Geo metadata + reverse geocoding:} resolve GPS EXIF to city/landmark. Enables a
  powerful hybrid --- a \emph{hard SQL geo filter} (``in Paris'') composed with \emph{soft visual
  ranking}, exactly the pre-filter pattern from \S6.
  \item \textbf{Landmark recognition} (``in front of the Eiffel Tower'') via embedding lookup
  against a landmark gallery --- another region-like signal.
  \item \textbf{Garment$\leftrightarrow$weather consistency prior:} reward agreement (raincoat
  $\leftrightarrow$ rain, parka $\leftrightarrow$ snow), capturing context a query \emph{implies}
  but does not state --- enabling \emph{``dressed for bad weather''}.
\end{enumerate}

Each is ``another axis + another weighted signal'', slotting into the existing fusion with no change
to the retrieval flow.

\subsection{Improving precision}

\begin{enumerate}
  \item \textbf{Build a labelled validation set} --- highest value. Without labels, tuning is
  guesswork; with a few hundred \tcode{(query, relevant\_id)} pairs we get \textbf{precision@k /
  nDCG} and every change becomes measurable. The harness already accepts \tcode{--labels}.
  \item \textbf{Learn the fusion:} logistic regression or \textbf{LambdaMART} over the four signals
  replaces hand-tuned constants and naturally learns \emph{query-dependent} weighting.
  \item \textbf{Contrastive fine-tuning on hard negatives:} LoRA-adapt FashionCLIP using
  \textbf{colour-swapped pairs} as explicit negatives, pushing binding into the \emph{encoder} ---
  attacking the root cause in \S1.
  \item \textbf{Stronger colour grounding:} ensemble the CLIP colour head with a classical
  \textbf{CIELAB} dominant-colour estimate per region (robust to lighting) --- should lift queries
  1, 3, 5.
  \item \textbf{Per-person binding} (\S5) --- the biggest correctness fix for crowded scenes.
  \item \textbf{Cross-encoder reranking} --- implemented behind a flag, to be enabled only once its
  uplift is \emph{measured}.
  \item \textbf{Query expansion} for style inference: expand \emph{``weekend casual''} into likely
  garments via a learned style$\rightarrow$garment prior.
\end{enumerate}

% ---------------------------------------------------------------- appendix
\section*{Appendix --- Reproducing}
\addcontentsline{toc}{section}{Appendix}

\begin{quote}\small\ttfamily
pip install -r requirements.txt\\
python main.py index \hfill \# Part A --- build FAISS + SQLite\\
python main.py query "a red tie and a white shirt" -k 10 \hfill \# Part B --- top-k retrieval\\
python main.py evaluate \hfill \# 5 official queries\\
python main.py compositional \hfill \# the colour-swap test of \S4.1\\
python webapp/app.py \hfill \# optional visual portal\\
pytest -q \hfill \# unit tests
\end{quote}

\textbf{Stack:} FashionCLIP (\tcode{patrickjohncyh/fashion-clip}), BLIP
(\tcode{blip-image-captioning-base}), MPNet (\tcode{all-mpnet-base-v2}), SegFormer
(\tcode{mattmdjaga/segformer\_b2\_clothes}), FAISS, SQLite. Hardware: one 4\,GB NVIDIA RTX 2050.

\vspace{4pt}
\hrule height 0.4pt
\vspace{4pt}
\begin{center}\small
Codebase: \href{\repourl}{\repourl}
\end{center}

\end{document}
